\documentclass[twocolumn]{iopjournal}

\usepackage{amsmath,amssymb,amsfonts}

\fancyhfoffset[L]{0mm}
\setlength{\marginparwidth}{0pt}
\setlength{\marginparsep}{0pt}

% ==========================================
% ZENODO PREPRINT HEADERS (ACTIVE)
% ==========================================
\renewcommand{\articletype}[1]{%
  {\vspace*{-8mm}\noindent \Large \sf Preprint}\par
  \vspace*{6mm}\noindent{\scriptsize \sf{\bfseries \MakeUppercase{#1}}}}

\fancyhead[L]{}
\fancyhead[R]{}

\hypersetup{pdftitle={The Two-Potential Energy of the Constitutive Law in Bimetric MOND}, pdfauthor={Stephen Nelson}, pdfsubject={Gravity Field Dynamics (GFD)}, pdfkeywords={BIMOND, QUMOND, AQUAL, MOND, bimetric gravity, Legendre duality, proof}}

\begin{document}

\articletype{Supplementary Material}

\title{The Two-Potential Energy of the Constitutive Law\\
$\mu(x) = x/(x+1)$ in Bimetric MOND}

\author{Stephen Nelson}

\affil{Field Dynamics Research}

\ead{stephen@fielddynamics.org}

\begin{abstract}
Milgrom's bimetric MOND (BIMOND) has two free choices in its
nonrelativistic limit: the ratio of the two Einstein--Hilbert
couplings $\alpha,\beta$ and an interaction function $\mathcal{M}$,
which is left unspecified apart from its two limits. We show that
the difference potential $\phi^* = \phi - \hat\phi$ obeys a Poisson
equation independent of $\mathcal{M}$ if and only if
$\alpha + \beta = 0$, and that in the tetrad form of the
Einstein--Hilbert term a frame on the second root of the tetrad
equation enters with the opposite sign, so equal couplings give
exactly this case. Requiring that BIMOND realise the constitutive
law $\mu(x) = x/(x+1)$ then fixes $\mathcal{M}$ uniquely:
$\mathcal{M} = 2x - 1 + 2b - b^2 + 2\ln b$, with
$b = 1/(1+x)$ and $x$ the total field in units of $a_0$. It meets
Milgrom's two conditions exactly. It is the Legendre dual of the
AQUAL kinetic function $F$, with the logarithms of the two cancelling
in the dual pairing, and its energy is strictly convex, with a
Hessian ratio that is the reciprocal of the AQUAL one.
\end{abstract}

\keywords{BIMOND, QUMOND, AQUAL, MOND, bimetric gravity, Legendre duality, proof}

% ----------------------------------------------------------------------

\section*{Setting}

The nonrelativistic limit of BIMOND~\cite{Milgrom2009,Milgrom2010twin}
has two potentials: $\phi$, felt by matter of density $\rho$, and
$\hat\phi$, felt by twin matter of density $\hat\rho$. With
$\phi^* = \phi - \hat\phi$ and $z = |\nabla\phi^*|^2/a_0^2$ the
Lagrangian density is
\begin{equation}
\begin{aligned}
  \mathcal{L} = &-\frac{1}{8\pi G}\Big[\beta|\nabla\phi|^2
     + \alpha|\nabla\hat\phi|^2 - a_0^2\,\mathcal{M}(z)\Big] \\
     &- \rho\,\phi + \hat\rho\,\hat\phi ,
\end{aligned}
  \label{eq:L}
\end{equation}
with the twin sign of Ref.~\cite{Milgrom2010twin}. Its
Euler--Lagrange equations are
\begin{equation}
\begin{aligned}
  &\nabla\cdot\big[\beta\nabla\phi - \mathcal{M}'\nabla\phi^*\big]
    = 4\pi G\rho, \\
  &\nabla\cdot\big[\alpha\nabla\hat\phi + \mathcal{M}'\nabla\phi^*\big]
    = -4\pi G\hat\rho ,
\end{aligned}
  \label{eq:EL}
\end{equation}
where $\mathcal{M}' = d\mathcal{M}/dz$. An additive constant in
$\mathcal{M}$ leaves Eq.~\eqref{eq:EL} unchanged, and we fix it by
$\mathcal{M}(0) = 0$. Throughout,
\begin{equation}
  \mu(x) = \frac{x}{x+1}.
  \label{eq:mu}
\end{equation}

% ----------------------------------------------------------------------

\section*{The Coupling Ratio}

\paragraph{Theorem 1.} \textit{Let $\alpha\beta \neq 0$ and let
$\mathcal{M}'$ be nonconstant. The potential $\phi^*$ obeys a linear
equation independent of $\mathcal{M}$ if and only if
$\alpha + \beta = 0$. With $\beta = 1$ the system is then}
\begin{equation}
\begin{aligned}
  \Delta\phi^* &= 4\pi G\,(\rho - \hat\rho), \\
  \Delta\phi &= 4\pi G\rho + \nabla\cdot\big(\mathcal{M}'\nabla\phi^*\big).
\end{aligned}
  \label{eq:qumond}
\end{equation}

\paragraph{Proof.} $\alpha$ times the first of Eqs.~\eqref{eq:EL}
minus $\beta$ times the second gives
\begin{equation}
  \nabla\cdot\Big[\big(\alpha\beta - (\alpha+\beta)\mathcal{M}'\big)
  \nabla\phi^*\Big] = 4\pi G(\alpha\rho + \beta\hat\rho),
  \label{eq:star}
\end{equation}
the unique combination, up to a factor, in which $\hat\phi$ cancels. Its
coefficient is a constant exactly when $(\alpha+\beta)\mathcal{M}'$
is, that is, when $\alpha + \beta = 0$. Setting $\alpha = -1$,
$\beta = 1$ in Eq.~\eqref{eq:star} and in the first of
Eqs.~\eqref{eq:EL} gives Eq.~\eqref{eq:qumond}. \hfill $\square$

\medskip
\noindent
A MOND theory needs a nonconstant $\mathcal{M}'$, since it must tend
to $0$ at high acceleration and diverge at low acceleration. For
$\alpha + \beta = 0$, Eq.~\eqref{eq:qumond} is the quasilinear
formulation QUMOND~\cite{Milgrom2010qumond}: two Poisson equations,
the second sourced by the first.

\paragraph{Lemma 1 (orientation).} \textit{Write the Einstein--Hilbert
term of a frame $e^a{}_\mu$ in its tetrad form,
$S[e] = (16\pi G)^{-1}\!\int \det(e)\,R[g(e)]\,d^4x$ with
$g = e^T\eta e$, the 4-form
$\epsilon_{abcd}\,e^a\wedge e^b\wedge R^{cd}$ up to normalisation.
For constant $\Lambda \in O(1,3)$, $S[\Lambda e] = \det(\Lambda)\,S[e]$.
In particular spatial parity $P$ gives $S[Pe] = -S[e]$.}

\paragraph{Proof.} $(\Lambda e)^T\eta(\Lambda e) = e^T\Lambda^T\eta
\Lambda e = e^T\eta e$, so $g$ and $R[g]$ are unchanged, while
$\det(\Lambda e) = \det(\Lambda)\det(e)$. Equivalently, in the
4-form, $\epsilon_{abcd}\Lambda^a{}_{a'}\Lambda^b{}_{b'}
\Lambda^c{}_{c'}\Lambda^d{}_{d'} = \det(\Lambda)\,\epsilon_{a'b'c'd'}$.
\hfill $\square$

\medskip
\noindent
The tetrad equation $g = e^T\eta e$ has two roots,
$\det e = \pm\sqrt{-g}$, exchanged by $P$~\cite{TwoRegions}. If the
second metric $\hat g$ is carried by a frame on the negative root
and both frames have the same coupling $(16\pi G)^{-1}$, then
$\det(\hat e)\,\hat R = -\sqrt{-\hat g}\,\hat R$, so in BIMOND's
metric notation $\beta = 1$ and $\alpha = -1$: the case of
Theorem~1.

% ----------------------------------------------------------------------

\section*{The Interaction Function}

\paragraph{Definition.} \textit{BIMOND realises $\mu$ if, for
$\hat\rho = 0$, its spherically symmetric fields coincide with those
of AQUAL~\cite{BekensteinMilgrom1984} with constitutive law $\mu$.}

\paragraph{Proposition 1.} \textit{Write $x = |\nabla\phi|/a_0$ for
the total field and $y = |\nabla\phi^*|/a_0$. BIMOND with
$\alpha + \beta = 0$ realises $\mu$ if and only if
$1 + \mathcal{M}'(y^2) = \nu(y)$, where $\nu(y) = x/y$ and
$y = x\,\mu(x)$. For Eq.~\eqref{eq:mu},}
\begin{equation}
  x(y) = \frac{y + \sqrt{y^2 + 4y}}{2},\qquad
  \nu(y) = \frac12 + \sqrt{\frac14 + \frac1y}.
  \label{eq:nu}
\end{equation}

\paragraph{Proof.} In spherical symmetry Gauss's law applied to
Eq.~\eqref{eq:qumond} gives $|\nabla\phi^*| = g_N$, the Newtonian
field, and $|\nabla\phi| = (1 + \mathcal{M}')\,g_N$. AQUAL gives
$\mu(|\nabla\phi|/a_0)\,|\nabla\phi| = g_N$. These agree for every
$g_N$ exactly when $1 + \mathcal{M}'$ is the inverse partner $\nu$.
Since $y = x^2/(1+x)$ increases from $0$ to $\infty$, the inverse is
unique and is Eq.~\eqref{eq:nu}. \hfill $\square$

\paragraph{Theorem 2.} \textit{The unique $\mathcal{M}$ with
$1 + \mathcal{M}'(y^2) = \nu(y)$ and $\mathcal{M}(0) = 0$ is, with
$b = 1/(1+x)$,}
\begin{equation}
  \boxed{\mathcal{M} = 2x - 1 + 2b - b^2 + 2\ln b ,}
  \label{eq:M}
\end{equation}
\textit{equivalently
$\mathcal{M} = x(2x^2+3x+2)/(x+1)^2 - 2\ln(1+x)$, or, in $y$,}
\begin{equation}
  \mathcal{M} = \tfrac12\sqrt y\,(y+2)\sqrt{y+4} - \tfrac12 y^2
     - 4\,\mathrm{arsinh}\big(\tfrac12\sqrt y\big).
  \label{eq:My}
\end{equation}
\textit{Its derivative is $\mathcal{M}'(z) = -\tfrac12 +
\sqrt{\tfrac14 + z^{-1/2}}$.}

\paragraph{Proof.} Two such functions differ by a constant, which
$\mathcal{M}(0) = 0$ removes. Since $x - y = x/(1+x) = \mu(x)$,
\begin{equation}
  \frac{d\mathcal{M}}{dy} = 2y\,\mathcal{M}'(y^2) = 2(x - y)
  = 2\,\mu(x).
  \label{eq:dM}
\end{equation}
With $dy/dx = x(x+2)/(x+1)^2$ and $t = s + 1$,
\begin{equation}
\begin{aligned}
  \mathcal{M} &= 2\int_0^x \frac{s^2(s+2)}{(s+1)^3}\,ds \\
  &= 2\int_1^{1+x}\Big(1 - \frac1t - \frac1{t^2} + \frac1{t^3}\Big)dt,
\end{aligned}
\end{equation}
which integrates to Eq.~\eqref{eq:M}. Eq.~\eqref{eq:My} has the same
derivative $2(x - y)$ in $y$ and vanishes at $y = 0$.
\hfill $\square$

\paragraph{Corollary 1 (the two limits).} \textit{As $z \to \infty$,
$\mathcal{M}'(z) \to 0$; as $z \to 0$,
$\mathcal{M}'(z) = z^{-1/4} - \tfrac12 + \tfrac18 z^{1/4} + O(z^{3/4})$.
Correspondingly}
\begin{equation}
\begin{aligned}
  \mathcal{M} &= \tfrac43 y^{3/2} - \tfrac12 y^2 + \tfrac1{10}y^{5/2}
     + O(y^3),\\
     &\qquad y \to 0,\\
  \mathcal{M} &= 2y - 2\ln y + 1 - \frac4y + \frac5{y^2} + O(y^{-3}),\\
     &\qquad y \to \infty .
\end{aligned}
  \label{eq:limits}
\end{equation}

\paragraph{Proof.} Expand Eq.~\eqref{eq:nu} and Eq.~\eqref{eq:My}.
\hfill $\square$

\medskip
\noindent
These are exactly the two conditions BIMOND imposes on
$\mathcal{M}$~\cite{Milgrom2009}: $\mathcal{M}' \to 0$ in the
Newtonian regime and $\mathcal{M}' \approx z^{-1/4}$ in the deep-MOND
regime, here with unit coefficient and no free normalisation.

% ----------------------------------------------------------------------

\section*{Legendre Duality and Convexity}

Let $g(x) = F(x^2) = x^2 - 2x + 2\ln(1+x)$ be the AQUAL kinetic
function of $\mu$~\cite{AqualProof}, and $Q(y) = y^2 + \mathcal{M}$
the kinetic function of Eq.~\eqref{eq:qumond}: the source of
$\Delta\phi$ is the divergence of $\nu(y)\nabla\phi^*$, and
$\nu(y)\,y = \tfrac12 Q'(y)$ by Eq.~\eqref{eq:dM}.

\paragraph{Theorem 3.} \textit{At $y = x\mu(x)$,}
\begin{equation}
\begin{aligned}
  &\boxed{g(x) + Q(y) = 2xy,}\\
  &g'(x) = 2y,\qquad Q'(y) = 2x,
\end{aligned}
  \label{eq:legendre}
\end{equation}
\textit{so $Q(y) = \max_{x \ge 0}\,[2xy - g(x)]$: $Q$ is the Legendre
transform of $g$. Written in $b$,
$g = x^2 - 2x - 2\ln b$ and $\mathcal{M} = 2x - 1 + 2b - b^2 + 2\ln b$:
the logarithms are equal and opposite and cancel in $g + \mathcal{M}$.
In spherical symmetry $|\nabla\hat\phi| = a_0\,\mu(x) < a_0$.}

\paragraph{Proof.} $g'(x) = 2x^2/(1+x) = 2y$ and, by
Eq.~\eqref{eq:dM}, $Q'(y) = 2y + 2(x - y) = 2x$. Then
$\frac{d}{dx}[g + Q - 2xy] = 2y + 2x\,y' - 2y - 2x\,y' = 0$ and all
three terms vanish at $x = 0$. Since $g'' > 0$ for $x > 0$, the
stationary point $g'(x) = 2y$ of $2xy - g(x)$ is its unique maximum.
Finally $|\nabla\hat\phi| = |\nabla\phi| - |\nabla\phi^*| =
a_0(x - y) = a_0\mu(x)$. \hfill $\square$

\paragraph{Corollary 2 (convexity and uniqueness).} \textit{The energy
$W(q) = Q(|q|)$, $q = \nabla\phi^*/a_0$, is strictly convex for
$q \neq 0$, with Hessian eigenvalues $Q'/y = 2\nu$ (twice) and
$Q'' = 2\,dx/dy = 2(x+1)^2/[x(x+2)]$ (once), whose ratio}
\begin{equation}
  \frac12 \le \frac{Q''}{Q'/y} = \frac{x+1}{x+2} < 1
  \label{eq:ratio}
\end{equation}
\textit{is the reciprocal of the AQUAL ratio $(x\mu)'/\mu =
(x+2)/(x+1)$. On a bounded domain with given boundary values,
Eq.~\eqref{eq:qumond} has at most one solution $(\phi^*, \phi)$.}

\paragraph{Proof.} For a radial function the Hessian has eigenvalue
$Q'/y$ across $\hat q$ and $Q''$ along it. $Q'/y = 2x/y = 2\nu$ and
$Q'' = 2/(dy/dx)$, both positive. Their ratio is
$(dx/dy)(y/x) = 1/[(x\mu)'/\mu]$, since $y = x\mu$; this is the
inverse Hessian of a Legendre pair. For uniqueness, the difference
of two solutions of each Poisson equation in Eq.~\eqref{eq:qumond}
is harmonic with zero boundary values, hence zero: first for
$\phi^*$, then for $\phi$, whose source is then fixed. The flux
$\nu(y)\,y = x(y)$ is continuous at $\nabla\phi^* = 0$, where
$x \simeq \sqrt y$, so the source $\nabla\cdot(\mathcal{M}'\nabla
\phi^*)$ is well defined weakly. \hfill $\square$

\medskip
\noindent
Where AQUAL degenerates at $\nabla\Phi = 0$ (both eigenvalues $\to 0$),
the two-potential energy is singular there ($2\nu \to \infty$); the
Legendre transform exchanges the two behaviours.

% ----------------------------------------------------------------------

\section*{Summary}

In nonrelativistic BIMOND the difference potential is Newtonian,
independently of the interaction, exactly when $\alpha + \beta = 0$,
and the tetrad form of the Einstein--Hilbert term gives that sign
for a frame on the second root. Requiring the constitutive law
$\mu(x) = x/(x+1)$ then fixes the interaction function uniquely as
$\mathcal{M} = 2x - 1 + 2b - b^2 + 2\ln b$, which meets BIMOND's
two limiting conditions with unit coefficient. It is the Legendre
dual of the AQUAL kinetic function, the logarithms of the two
cancelling in the pairing, and its energy is strictly convex with
the reciprocal Hessian ratio.

% ----------------------------------------------------------------------

\begin{thebibliography}{9}

\bibitem{Milgrom2009}
M.~Milgrom,
\textit{Bimetric MOND gravity},
Phys.\ Rev.\ D \textbf{80} (2009) 123536, arXiv:0912.0790.

\bibitem{Milgrom2010twin}
M.~Milgrom,
\textit{Matter and twin matter in bimetric MOND},
arXiv:1001.4444 (2010).

\bibitem{Milgrom2010qumond}
M.~Milgrom,
\textit{Quasi-linear formulation of MOND},
Mon.\ Not.\ R.\ Astron.\ Soc.\ \textbf{403} (2010) 886,
arXiv:0911.5464.

\bibitem{BekensteinMilgrom1984}
J.~Bekenstein and M.~Milgrom,
\textit{Does the missing mass problem signal the breakdown of
Newtonian gravity?},
Astrophys.\ J.\ \textbf{286} (1984) 7.

\bibitem{AqualProof}
S.~Nelson,
\textit{The energy functional of the constitutive law
$\mu(x) = x/(x+1)$}, Field Dynamics Research (2026).

\bibitem{TwoRegions}
S.~Nelson,
\textit{The tetrad scheme and the two bounded regions},
Field Dynamics Research (2026).

\end{thebibliography}

\end{document}
